\documentclass[11pt]{article}

\usepackage[utf8]{inputenc}
\usepackage[T1]{fontenc}
\usepackage{lmodern}
\usepackage{geometry}
\usepackage{amsmath,amssymb,amsfonts,amsthm,mathtools}
\usepackage{enumitem}
\usepackage{microtype}
\usepackage{hyperref}

\geometry{margin=1in}
\hypersetup{colorlinks=true, linkcolor=black, urlcolor=blue, citecolor=black}
\setlist[enumerate]{topsep=0.35em,itemsep=0.3em}
\setlist[itemize]{topsep=0.3em,itemsep=0.25em}

\newcommand{\Aether}{\AE{}ther}
\newcommand{\Aflow}{\AE{}ther-flow}
\newcommand{\TheoryProgram}{\Aether{} / \Aflow{} framework}
\newcommand{\Lang}{\mathcal{L}_{0}}
\newcommand{\LangPlus}{\mathcal{L}_{0}^{+H}}
\newcommand{\LedgerFam}{\mathfrak{H}}
\newcommand{\PiH}{\Pi_H}
\newcommand{\AdH}{\mathsf{Ad}_{H}}
\newcommand{\BlindH}{\mathsf{Blind}_{H}}
\newcommand{\GenH}{\mathsf{Gen}_{H}}
\newcommand{\EnvH}{\mathsf{Env}_{H}}
\newcommand{\EqSrc}{\mathsf{Eq}_{\mathrm{src}}}
\newcommand{\RetainH}{\mathsf{Retain}_{H}}
\newcommand{\NoTargetH}{\mathsf{NoTarget}_{H}}
\newcommand{\AuditH}{\mathsf{Audit}_{H}}
\newcommand{\Current}{\mathsf{Current}}
\newcommand{\Neg}{\mathsf{Neg}}
\newcommand{\Expand}{\mathsf{Expand}}
\newcommand{\Unres}{\mathsf{Unresolved}}
\newcommand{\Pass}{\mathsf{Pass}}
\newcommand{\Fail}{\mathsf{Fail}}
\newcommand{\Block}{\mathsf{Block}}

\newtheorem{test}{Refutation test}
\newtheorem{finding}{Refutation finding}
\newtheorem{variation}{Admissible variation}
\newtheorem{blocker}{Promotion blocker}
\newtheorem{verdict}{Local verdict}
\newtheorem{nextburden}{Next burden}

\title{\textbf{Localization Source-Basis Axiom Selector-Domain Ledger-Generation Ontology-Expansion Proposal Conservativity and Benchmark-Independence Refutation}\\[0.35em]
\large Variation Test for Provenance, Admissibility, Blinding, Source Equivalence, Retention, and Target Exclusion}
\author{Refuter AgentJob AJ-RT-20260613-016-001}
\date{June 13, 2026}

\begin{document}

\maketitle

\begin{abstract}
This draft/control artifact stress-tests the audited noncanonical
ledger-generation ontology-expansion proposal for conservativity and
benchmark independence. The prior Smuggling Auditor result is preserved: the
proposal is target-clean as a proposal and does not directly claim protected
authority. The present result is local and negative for candidate-support
use. Conservativity over \(\Lang\) is not proved, and benchmark independence
is not established, because admissible variation of provenance packets,
admissibility rules, blinding records, source equivalence, retention schemas,
and target-exclusion grammars can change the generated ledger family or hide
adverse alternatives. The proposal remains a noncanonical research object.
Candidate reconstruction, canonical ontology edit, benchmark promotion, Gate
Chair review, and completed-derivation language remain blocked.
\end{abstract}

\section{Scope and Refutation Standard}

This artifact is a draft/control output for task \texttt{RT-20260613-016}.
It does not edit canonical ontology TeX, adopt \(\LangPlus\) as ontology,
authorize candidate reconstruction, issue a Gate Chair verdict, promote
benchmark status, or reject the broader \TheoryProgram{}.

The stress-test target is the audited proposal
\[
    \LangPlus =
    \Lang \cup
    \{\AdH,\BlindH,\GenH,\EqSrc,\RetainH,\NoTargetH,\AuditH\}.
\]
The proposal passed a narrow smuggling audit as a proposal. That audit did
not prove that the extension is conservative, that \(\GenH\) is complete,
that \(\EqSrc\) preserves adverse behavior, that \(\RetainH\) retains all
non-equivalent ledgers with relevant outputs and statuses, or that
\(\NoTargetH\) is a mechanical source-syntax grammar.

\begin{test}[Conservativity and benchmark independence]
The proposal can become candidate-support evidence only if the following
two local tests survive:
\begin{enumerate}
    \item \textbf{Conservativity}: adding \(\LangPlus\) must not imply new
    candidate reconstruction, benchmark success, target metric recovery,
    target-locality recovery, target observer protocol, target proper-time
    normalization, target-favorable rank, or canonical ontology claims over
    pre-existing \(\Lang\)-claims.
    \item \textbf{Benchmark independence}: admissible source-side variation
    of provenance packets, admissibility rules, blinding records,
    source-equivalence grammars, retention schemas, and target-exclusion
    grammars must not select a favorable ledger or suppress adverse ledgers
    except by declared source theorem.
\end{enumerate}
\end{test}

\section{Provenance-Packet Variation}

\begin{finding}[The provenance packet is not source-forced]
The proposed input
\[
\Pi_H=(\mathcal{R}_{0},\mathcal{C}_{0},\mathcal{O}_{0},
\mathcal{Q}_{0},\mathcal{B}_{0})
\]
names source records, source closure rules, source operation seeds, source
quotient seeds, and pre-benchmark blinding records. Naming the packet does
not prove that this packet is the unique or complete source-generated packet
inside \(\Lang\).
\end{finding}

\begin{variation}[Packet breadth and seed variation]
Two source-clean packets can differ in closure breadth, quotient refinement,
operation seeds, or blinding records while both remain free of target metric
behavior and benchmark-success predicates. One packet can generate a
family in which a favorable ledger appears distinguished; another can
generate additional adverse ledgers or require protected expansion.
\end{variation}

\begin{blocker}[Packet dependence blocks benchmark independence]
Until a source theorem fixes the allowed packet class or retains all
non-equivalent packet-generated families, the proposal is sensitive to
packet choice. A favorable result under one packet cannot support
candidate reconstruction.
\end{blocker}

\section{Admissibility-Rule Variation}

\begin{finding}[Admissibility breadth is a live theorem burden]
\(\AdH(\Pi,H')\) classifies admissible ledgers, but the proposal does not
give a theorem that fixes the breadth of admissibility under narrower,
broader, and incomparable source rules.
\end{finding}

\begin{variation}[Narrow and broad admissibility]
A narrow admissibility rule can exclude adverse ledgers as source-invalid.
A broader rule can include ledgers with different presentation boundaries,
fronts, quotient spectra, transport reports, or status rows. Both rules can
be written without direct target vocabulary.
\end{variation}

\begin{blocker}[Admissibility variation can reintroduce selection]
If \(\AdH\) is chosen so that only candidate-friendly ledgers remain, the
proposal has reintroduced selection through admissibility. If \(\AdH\) is
broad enough to include adverse alternatives, candidate support remains
blocked unless complete retention and source equivalence are proved.
\end{blocker}

\section{Blinding-Record Variation}

\begin{finding}[Blinding is not yet a mechanical independence theorem]
\(\BlindH(\Pi,H')\) requires construction without benchmark outcome
inspection. The proposal does not prove that the blinding record rules are
themselves invariant under admissible record formats, hashes, time locks,
pre-registration evidence, or source replay procedures.
\end{finding}

\begin{variation}[Blinding formats with different evidential reach]
One blinding record may seal only direct target-output inspection. Another
may also seal downstream candidate-performance heuristics, target-favorable
rank, or proxy outcome labels. A third may preserve source replay while
leaving relevance choices underdetermined.
\end{variation}

\begin{blocker}[Blinding variation leaves proxy-selection risk]
Without a theorem that \(\BlindH\) excludes both direct target inspection and
source-side proxy selection, the proposal cannot prove benchmark
independence. Passing a narrow blinding record is not enough.
\end{blocker}

\section{Source-Equivalence Variation}

\begin{finding}[Equivalence after generation is necessary but not sufficient]
The proposal correctly places \(\EqSrc\) after generation. The unresolved
question is whether \(\EqSrc\) is an equivalence relation whose quotient
classes preserve all adverse outputs and statuses relevant to negative
controls.
\end{finding}

\begin{variation}[Coarse and fine quotient grammars]
A coarse source-equivalence grammar can collapse ledgers that differ in
adverse source behavior. A fine grammar can split ledgers into multiple
non-equivalent alternatives. Both can be described as source relabeling
unless the grammar states invariants and failure conditions.
\end{variation}

\begin{blocker}[Quotient variation can erase adverse ledgers]
If \(\EqSrc\) erases adverse behavior before retention, the proposal fails
benchmark independence. If \(\EqSrc\) is too fine or underspecified, the
generated family remains non-unique and candidate support remains blocked.
\end{blocker}

\section{Retention-Schema Variation}

\begin{finding}[Retention completeness is not proved]
\(\RetainH\) states that non-equivalent generated ledgers remain retained as
negative-control alternatives. It does not prove that all such ledgers are
enumerated or schema-retained with their outputs and statuses.
\end{finding}

\begin{variation}[Representative retention versus full retention]
One retention schema can retain only representative ledgers or only ledger
identifiers. Another can retain full source outputs, quotient status,
transport reports, holonomy reports, and \(D_F\) status rows. A third can
retain incomplete placeholders marked unresolved.
\end{variation}

\begin{blocker}[Incomplete retention permits favorable projection]
Candidate support requires adverse alternatives to remain visible at the
same evidential grain as favorable alternatives. If \(\RetainH\) retains
less information for adverse ledgers than for favorable ledgers, the
proposal permits projection rather than benchmark independence.
\end{blocker}

\section{Target-Exclusion Grammar Variation}

\begin{finding}[Target exclusion is not yet source syntax]
\(\NoTargetH\) is useful as a guard, but the proposal has not converted it
into a mechanical source-syntax grammar. A list of forbidden target channels
does not itself decide which source expressions, labels, proxy ranks, or
relevance choices are excluded.
\end{finding}

\begin{variation}[Syntactic and semantic exclusion]
A syntactic \(\NoTargetH\) may block explicit target metric or benchmark
labels while leaving source-side proxies for target success. A broader
semantic \(\NoTargetH\) may require authority not present in \(\Lang\)
because deciding proxy equivalence can itself require target knowledge.
\end{variation}

\begin{blocker}[Target-exclusion variation forces controlled pause]
Until \(\NoTargetH\) is specified as source syntax with failure conditions,
the stress test cannot distinguish source-clean selection from disguised
target-performance selection. The safe result is controlled pause.
\end{blocker}

\section{Stress-Test Ledger}

\begin{center}
\begin{tabular}{p{0.18\linewidth}p{0.36\linewidth}p{0.31\linewidth}}
\textbf{Surface} & \textbf{Admissible variation} & \textbf{Local result} \\
\hline
\(\Pi_H\) & records, closure rules, operation seeds, quotient seeds, blinding records & packet dependence not source-forced \\
\(\AdH\) & narrow, broad, or incomparable admissibility & favorable selection can reappear through admissibility breadth \\
\(\BlindH\) & hashes, time locks, replay evidence, proxy-outcome scope & proxy-selection exclusion not proved \\
\(\EqSrc\) & coarse or fine quotient grammar and invariants & adverse ledgers can be erased or remain non-unique \\
\(\RetainH\) & representative, full, or placeholder retention & complete adverse retention not proved \\
\(\NoTargetH\) & syntactic exclusion, semantic proxy exclusion, failure conditions & target-clean source grammar not mechanized \\
\end{tabular}
\end{center}

\section{Local Verdict}

\begin{verdict}[Conservativity and benchmark independence are not established]
The audited proposal remains a coherent noncanonical object for further
scrutiny, but it does not pass the present Refuter test as candidate-support
evidence. Conservativity over \(\Lang\) is not proved because \(\LangPlus\)
adds generator, admissibility, equivalence, retention, and target-exclusion
vocabulary whose consequences over existing source claims are not bounded by
theorem. Benchmark independence is not proved because admissible variation
can change the generated family, quotient adverse ledgers, incompletely
retain alternatives, or leave source-side proxy selection unresolved.
\end{verdict}

\begin{verdict}[Candidate-support use remains blocked]
The negative result is local. It does not reject the \TheoryProgram{} and
does not reject future ontology-expansion repair work. It blocks candidate
reconstruction, canonical ontology edit, benchmark promotion, Gate Chair
review, and completed-derivation language from the current proposal.
\end{verdict}

\section{Next Burden}

\begin{nextburden}[Mechanization repair target or controlled pause]
The logical next role is not Candidate Constructor and not Gate Chair
review. A bounded Ontology Formalizer packet may isolate the least broad
mechanization burden among \(\NoTargetH\), \(\EqSrc\), \(\RetainH\), and
\(\GenH\): either specify a source-syntax theorem target with explicit
failure conditions, or preserve the branch as controlled-pause evidence.
Any such packet must remain noncanonical and must not authorize candidate
reconstruction, canonical ontology edit, benchmark promotion, or Gate Chair
review.
\end{nextburden}

\section*{References}

Ricciardi, A. S. (2026, June 13). \emph{Localization source-basis axiom
selector-domain ledger-generation ontology-expansion proposal smuggling
audit: Target-import, authority-laundering, quotient-laundering, retention,
and status review} [Unpublished manuscript]. The Aether-Flow Interpretation
of Relativity Research Project,
\texttt{research\_control/tasks/RT-20260613-015/artifacts/31\_LOCALIZATION\_SOURCE\_BASIS\_AXIOM\_SELECTOR\_DOMAIN\_LEDGER\_GENERATION\_ONTOLOGY\_EXPANSION\_PROPOSAL\_SMUGGLING\_AUDIT.tex}.

Ricciardi, A. S. (2026, June 13). \emph{Localization source-basis axiom
selector-domain ledger-generation ontology-expansion proposal: Noncanonical
source-language authority candidate for family-first ledger provenance}
[Unpublished manuscript]. The Aether-Flow Interpretation of Relativity
Research Project,
\texttt{research\_control/tasks/RT-20260613-010/artifacts/30\_LOCALIZATION\_SOURCE\_BASIS\_AXIOM\_SELECTOR\_DOMAIN\_LEDGER\_GENERATION\_ONTOLOGY\_EXPANSION\_PROPOSAL.tex}.

Ricciardi, A. S. (2026, June 13). \emph{Localization source-basis axiom
selector-domain ledger-provenance theorem attempt or controlled pause:
Single-surface current-language test for Pi\_H, H, and admissible ledger
retention} [Unpublished manuscript]. The Aether-Flow Interpretation of
Relativity Research Project,
\texttt{research\_control/tasks/RT-20260613-009/artifacts/29\_LOCALIZATION\_SOURCE\_BASIS\_AXIOM\_SELECTOR\_DOMAIN\_LEDGER\_PROVENANCE\_THEOREM\_ATTEMPT\_OR\_CONTROLLED\_PAUSE.tex}.

Ricciardi, A. S. (2026, June 13). \emph{Localization source-basis axiom
selector-domain source-forcing discharge protocol benchmark-independence
refutation: Stability test for four-cell ledger classification under
source-side variation} [Unpublished manuscript]. The Aether-Flow
Interpretation of Relativity Research Project,
\texttt{research\_control/tasks/RT-20260613-008/artifacts/28\_LOCALIZATION\_SOURCE\_BASIS\_AXIOM\_SELECTOR\_DOMAIN\_SOURCE\_FORCING\_DISCHARGE\_PROTOCOL\_BENCHMARK\_INDEPENDENCE\_REFUTATION.tex}.

\end{document}
